\chapter[Genera Studied]{List of Genera Studied Arranged by Their Plant Families}

At the beginning it was thought that there might be sufficient similarity in
germination behavior within families that the data could be organized by families. In
fact there proved to be not only much diversity within families, but sometimes within
genera. Nevertheless, some general comments about the larger families may be
helpful. The following list of families studied also provides an opportunity for listing the
genera in each family for which data was recorded.

\begin{description}
\item[Acanthaceae] Ruellia
\item[Aceraceae] Acer
\item[Aizoaceae] Delosperma, Mesembryanthemum
\item[Alismaceae] Alisma
\item[Amaryllidaceae] Alstroemeria, Galanthus, Habranthus, Hymenocallis,
Hypoxis, Ixolirion, Leucojum, Lycoris, Narcissus, Rhodolirion, Rhodophiala, Ungernia,
Zephyranthes. Many species germinate only at 40. Photorequirements were never
found.
\item[Ameranthaceae] Gomphrena
\item[Anacardiaceae] Rhus
\item[Annonaceae] Asimina
\item[Apiaceae (Umbelliferae)] Aciphylla, Anethum, Angelica, Anisotome,
Anthriscus, Astrantia, Athamanta, Bupleurum, Coriandrum, Cymopteris, Daucus,
Erigenia, Eryngium, Ferula, Heracleum, Levisticium, Lornatium, Meum, Myrrhis,
Pimpinella, Pseudotaenidia, Selinus, Zizia. Germinatons were complex with many
germinating at 40, some requiring light and being stimulated by GA-3, and oscillating
temperatures being required for Myrrhis.
\item[Apocyanaceae] Adenum, Amsonia, Nerium
\item[Aquifoliaceae] Ilex
\item[Araceae] Arisaema, Arum, Eminium, Lysichitum, Orontium, Pinellia,
Symplocarpus
\item[Araliaceae] Aralia, Panax
\item[Aristolochiaceae] Aristolochia, Asarum
\item[Asclepiadaceae] Asciepias, Matelea, Periploca
\item[Asteraceae (Compositae)] Achillea, Ajania, Anacyclus, Anaphalis,
Andryala, Antennaria, Anthemis, Arnica, Artemisia, Aster, Asteromoa, Balsamorhiza,
Carlina, Carthamus, Celmisia, Centaurea, Chaenactis, Chamaechaenactis,
Cremanthodium, Chrysanthemum, Chrysopsis, Cichorum, Coreopsis, Dimorphotheca,
Doronicum, Enceliopsis, Erigeron, Eriophyllum, Eupatorium, Gaillardia, Gutierrtzia,
Haplopappus, Hehianthus, Hehichrysum, Heterotheca, Hippolytica, Hymenoxys, Inula,
Leontopodium, Liatris, Ligularia, Lygodesmia, Machaeranthera, Marshallia,
Melampodium, Mutisia, Pachystegia, Psychrogeton, Rudbeckia, Saussurea, Senecio,
Solidago, Stokesia,Tanacetopsis, Tanacetum, Taraxacum, Townsendia, Tussilago,
Vemonia, Waidheimia, Wyethia, Xeranthemum, Xylanthemum. These are
predominantly D-70 germinators. One species of Aster, two Eupatorium, and one
Vernonia were 40-70 germinators. Centaurea has an interesting physical mechanism
for delaying germination.

\item[Balsaminaceae] Impatiens
\item[Begoniaceae] Begonia
\item[Berberidaceae] Berberis, Bongardia, Caulophyllum, Epimedum,
Gymnospermium, Jeffersonia, Leontice, Mahonia, Nandina, Podophyllum: These
have complex and multicycle germinations. Photoeffects, oscillating temperature
effects, and the common D-70 pattern were all absent.
\item[Betulaceae] Alhus, Betula, Corylus, Ostrya
\item[Bignoniaceae] Argylia, Campsis, Catalpa, Chilopsis, Incarvillea,
Jacaranda: Photorequirements were found in Campsis and Catalpa.
\item[Bixaceae] Bixa
\item[Boraginaceae] Anchusa, Arnebia, Borago, Cryptantha, Cynoglossum,
Echioides, Eritrichum, Lindelofia, Lithospermum, Mertensia, Myosotis, Onosma,
Pseudomertensia, Pulmonaria: While many are simple 0-70 germinators,
photorequirements and oscillating temperature requirements are present.
\item[Brassicaceae (Cruciferae)]. Aethionema, Alyssoides, Alyssum,
Arabis, Aubrieta, Brassica, Cheiranthus, Cochlearia, Degenia, Dentaria, Draba,
Erysimum, Eunomia, Fibigia, Hesperis, Hutchinsia, Iberis, Lepidium, Lesquerella,
Lunaria, Mathiola, Nasturtium, Parrya, Petrocallis, Physaria, Ptilot rich um,Stanleya,
Thlaspi: All are 0-70 germinators except Dentaria which requires oscillating
temperatures and germinates in a hypogeal manner whereas all other genera are
epigeal germinators.
\item[Butomaceae] Butomus
\item[Buxaceae] Sarcococca
\item[Cactaceae] Coryphantha, Echinocereus, Ferocactus, Neobessya,
Opuntia, Pediocactus, Rebutia, Scierocactus: Some species germinate readily at 70.
However, the discovery that E. pectinatus hybrids required gibberelins for germination
means that GA-3 treatment needs to be tried on the reluctant germinators.
\item[Calycanthaceae] Calycanthus
\item[Campanulaceae] Adenophora, Campanula, Codonopsis, Cyananthus,
Edraianthus, Jasione, Michauxia, Ostrowskia, Physoplexis, Phyteuma, Platycodon,
Symphyandra, Trachelium: Most are 0-70 germinators, but photorequirements are
found in some Campanula, and a GA-3 requirement was found in a few.
\item[Cannaceae] Canna
\item[Capparidaceae] Capparis, Cleome
\item[Caprifoliaceae] Kolkwitzia, Lonicera, Sambucus, Symphoricarpos,
Viburnum, Weigelia: These exhibit a variety of germination patterns with many
germinating exclusively at 40 whereas 40 is fatal to Kolkowitzia. Some germinations
are multistep. Photoeffects were not found.
\item[Caricaceae] Carica: Light was required.
\item[Caryophyllaceae] Acanthophyllum, Arenaria, Cerastium, Dianthus,
Gypsophila, Lychnis, Melandrium, Minuartia, Paronychia, Petrocoptis, Saponaria,
Silene: All are D-70 germinators except for Silene where photoeffects and other
complexities were found.
\item[Celastraceae] Celastrus, Euonymus: Some Euonymus require
oscillating temperatures, extended multicycles, and washing the seeds with.
detergents.
\item[Cercidiphyllaceae] Cercidiphyllum.
\item[Chenopodiaceae] Chenopodium, Morocarpus
\item[Chloranthaceae] Chloranthus
\item[Cistaceae] Cistus, Fumana, Relianthemum, Hudsonia
\item[Clethraceae] Clethra
\item[Commelinaceae] Commelina, Tradescantia
\item[Convolvulaceae] Convolvulus, lpomoea
\item[Coriariaceae] Coriaria
\item[Cornaceae] Aucuba, Cornus: Cornus generally require oscillating
temperatures.
\item[Crassulaceae] Chiastophyllum, Dudleya, Kalanchoe, Rhodiola, Sedum,
Sempervivum: Light and GA-3 requirements were found.
\item[Cucurbitaceae] Citrullus, Cucumis, Cucurbita, Echinocystis. Sicyos,
Trichosanthes: The temperate zone species require oscillating temperatures. The
germination behaviors are complex and interesting.
\item[Cupressaceae] Umbilicus
\item[Diapensiaceae] Pyxidanthera, Shortia: Shortia requres light and exogenous chemicals.
\item[Dilleniaceae] Actinidia
\item[Dioscoreaceae] Dioscorea, Tamus
\item[Dipsaceae] Cephalaria, Dipsacus, Jurinella, Scabiosa: Dipsacus requires light.
\item[Droseraceae] Dionea
\item[Ebenaceae] Diospyros
\item[Elaeagnaceae] Eleaganus, Hippophae
\item[Ephedraceae] Ephedra
\item[Ericaceae] Arbutus, Arctostaphylos, Bruckenthalia, Cassiope, Enkianthus,
Epigea, Gaultheria, Kalmia, Ledum, Leiophyllum, Leucothoe, Lyonia, Oxydendrum
Pernettya, Phyllodoce, Pieris, Rhododendron, Vaccinium: The majority more or less
require light.
\item[Euphorbiaceae] Euphorbia, Securinega.
\item[Fabaceae (Leguminosae)]. Acacia, Albizzia, Anthyllis, Astragalus,
Baptisia, Bauhinia, Cercis, Chamaecytisus, Cladastris, Clianthus, Colutea, Coronilla,
Cytisus, Dalea, Dorycnium, Genista, Gleditsia, Gymnocladus, Hardenburgia,
Hedysarum, Indigofera, Labumum, Lathyrus, Lespedeza, Lotus, Lupinus, Oxytropis,
Piptanthus, Robinia, Sophora, Tephrosia, Thermopsis, Trifolium, Vicia, Wisteria:
Impervious seed coats are present in most species.
\item[Fagaceae] Quercus
\item[Fumariaceae] Adlumia, Corydalis, Dicentra: These have some of the
most complex and extended germination patterns in the plant kingdom.
\item[Gentianaceae] Blackstonia, Centaurium, Eustoma, Frasera, Gentiana,
Gentianella, Gentianopsis, Lisianthus, Lomatogonum, Sabatia, Swertia:
Photorequirements and GA-3 requirements are present in some genera.
\item[Geraniaceae] Geranium
\item[Gesneriaceae] Briggsia, 1-laberlea, Jankae, Opithandra, Ramonda: All
are D-70 germinators. The pots of seedlings must be kept in high humidity for at least
a year as the seedlings are sensitive to dessication. This is readily accomplished by
keeping the pots in polyethylene bags which are loosely tied at the top with a wire.
\item[Gingkoaceae] Ginkgo
\item[Globulariaceae] Globularia
\item[Halorgidaceae] Gunnera
\item[Hamamelidaceae] Corylopsis, Hamamelis, Liquidambar, Parrotia,
Parrotiopsis: Germinations are usually multicycle.
\item[Hippocastanaceae] Aesculus
\item[Hydrangaceae] Deinanthe
\item[Hydrophyllaceae] Nemophila, Phacelia, Romanzoffia: Oscillating
temperatures are required in some species.
\item[Hypericaceae (Clusiaceae)]. Hypericurn: Light is usually required.
\item[Iridaceae] Acidanthera, Belamcanda, Crocosmia, Crocus, Dierama,
Freesia, Gladiolus, Hermodactylus, Iris, Lapeirousia, Melasphaerula, Moraea,
Nemastylis, Romulea, Sisyrinchium, Tigridia: A large variety of behaviors are found
including photorequirements and extended multicycle germinations.
\item[Juglandaceae] Juglans
\item[Lamiaceae] Aga\-stache, Ajuga, Ballota, Col\-lin\-sonia, Draco\-cephalum,
Eremo\-stachys, Horminum, Hymenocrater, Hyssopus, Lallemantia, Lamium,
Lavandula, Majorana, Mentha, Monarda, Monardella, Nepeta, Rosmarinus, Ocimum,
Origanum, Pelkovia, Petroselinum, Phlomis,Physostegia, Prunella, Salvia, Satureia,
Scutellaria, Sideritis, Stachys, Teucrium, Thymus, Trichostema, Ziziphora:
Photorequirements are found in several of the genera.
\item[Lardizabalaceae] Decaisnea
\item[Lauraceae] Lindera, Sassafras
\item[Lentibulariaceae] Pinguicula
\item[Liliaceae] Agave, Allium, Androcymbium, Androstephium, Anthericum,
Arthropodium, Asphodeline, Asphodelus, Asparagus, Bellevalia, Brimeura, Brodiaea,
Bulbinella, Bulbocodium, Calochortus, Caloscordum, Camassia, Cardiocrinum,
Chionodoxa, Chiorogalum, Clintonia, Colchicum, Convallaria, Disporum,Eremurus,
Erythronium, Fritillaria, Gagea, Galtonia, Hastingsia, Helonias, Hemerocallis, Hosta,
Hyacinthus, Ixolirion, Kniphofia, Korolkowia, Leucocrinum, Lilium, Liriope, Lloydia,
Maianthemum, Manfreda (Agave), Merendera, MuscariNothoscordum, Ophiopogon,
Ornithogalum, Paradisea, Paris, Polygonatum, Puschkinia, Rhinopetalum,
Sandersonia, Schoenolirion, Scilla, Scoliopus, Smilacina, Smilax, Strangweia,
Theropogon, Tricyrtis, Trillidium, Trillium, Triteleia, Tulipa, Uiularia, Veratrum,
Xeronema, Xerophyllum, Yucca, Zygadenus: Convallaria and Ophiopogon provided
the two clearest examples of blockage of germination by light. Tricyrtis generally
require light. Many species germinate only at 40. Multistep and multicycle
germinations are common. Oscillating temperature requirements are also found.
Every conceivable germination pattern is present.
\item[Limnanthaceae] Limnanthes
\item[Linaceae] Linum
\item[Loasaceae] Cajophora, Loasa, Mentzelia
\item[Lobeliaceae] Lobelia
\item[Loganiaceae] Buddleia, Spigelia
\item[Loranthaceae] Phoradendron
\item[Lythraceae] Cuphea, Lythrum
\item[Magnoliaceae] Magnolia
\item[Malesherbiaceae] Malesherbia
\item[Malvaceae] Abutilon, Hibiscus, Illiamna, Malva, Sidalcea, Sphaeralcea:
Impervious seed coats were found in some species.
\item[Melastomaceae] Osbeckia, Rhexia: Light was more or less required in
all species.
\item[Mesembrvanthemaceae] Delosperma
\item[Moraceae] Maclura, Morus: Germination is rapid if the inhibitors in the
fruit are removed.
\item[Myricaceae] Myrica
\item[Myrtaceae] Callistemon, Calothamnus, Eucalyptus, Kunzea, Melaleuca
\item[Nyctaginaceae] Abronia, Mirabilis
\item[Nymphaeaceae] Nelumbo, Nymphaea
\item[Nvssaceae] Nyssa
\item[Oleaceae] Abeliophyllum, Chionanthus, Forsythia, Fraxinus, Jasminum,
Ligustrum, Menodora, Syringa: Multistep and multicycle germinations were found.
\item[Onagraceae] Circaea, Clarkia, Epilobium, Fuschia Gaura, Oenothera,
Zauschneria: Photorequirements were found in a few species.
\item[Orchidaceae] See Chapter 21
\item[Oxalidaceae] Oxalis
\item[Palmaceae] Phoenix, Sabal, Washingtonia
\item[Papaveraceae] Argemone, Chelidonium, Glaucium, Macleaya,
Meconopsis, Papaver, Romneya, Sanguinaria, Stylophorum: A variety of patterns
were found including photorequirements and multicycle and multistep patterns.
\item[Parnassiaceae] Parnassia
\item[Passifloraceae] Passiflora
\item[Phytolaccaceae] Phytolacca
\item[Pinaceae] Abies, Cedrus, Cunninghamia, Cupressus, Juniperus, Larix,
Libocedrus, Picea, Pinus, Pseudotsuga, Sciadopitys, Sequoia, Sequoiadendron,
Taxodium, Thuja, Tsuga
\item[Platanaceae] Platanus
\item[Plumbaginaceae] Acantholimon, Armeria, Dictyolimon, Goniolimon,
Ikonnikovia, Limonium, Popoviolimon: These are generally D-70-germinators.
\item[Poaceae (Graminae)] Andropogon, Bouteloua, Briza, Bromus,
Calamagrostis, Eleusine, Festuca, Hystrix, Koeleria, Lagarus, Melica, Miscanthus,
Panicum, Pennisetum, Scleria, Stipa: See Chapter 22.
\item[Polemoniaceae] Collomia, Gilia, Ipomopsis, Leptodactylon, Leptosiphon,
Linanthrastum, Linanthus, Phlox, Polemonium: Germinations were varied and
complex. It is of much value to know the exact pattern.
\item[Polygalaceae] Polygala
\item[Polygonaceae] Eriogonum, Polygonum, Rheum, Rumex
\item[Portulacaceae] Calandrinia, Calyptridium, Claytonia, Lewisia, Montia,
Portulaca, Spraguea, Talinum: A variety of patterns were found including
photorequirements and low temperatures for germination.
\item[Primulaceae] Androsace, Cortusa, Cyclamen, Dionysia, Dodecatheon,
Douglasia, Primula, Soldanella: It is of much value to know the exact pattern as low
temperature germination, complex photoeffects, multicycle germinations, and
sensitivity to dry storage were all present.
\item[Proteaceae] Dryandra
\item[Puniceae] Punica
\item[Pyrolaceae] Pyrola
\item[Ranunculaceae] Aconitum, Actaea, Adonis, Anemone, Anemonastrum,
Anemonella, Anemonopsis, Aquilegia, Callianthemum, Caltha, Cimiifuga, Clematis,
Delphinium, Eranthis, Glaucidium, Helleborus, Hepatica, Hydrastis, Isopyrum,
Paeonia, Paraquilegia, Ranunculus, Thalictrum, Trollius: Low temperature, delayed,
and multicycle germination patterns were all present. Also found were
photorequirements, oscillating temperature requirements, and impervious seed coats.
Is it any wonder that this family has long been regarded as difficult to germinate. They
are really not difficult if the exact pattern is known.
\item[Rhamnaceae] Ceanothus, Rhamnus
\item[Rosaceae] Acaena, Alchemilla, Amelanchier, Armenaica, Aruncus,
Cercocarpus, Chaenomeies, Cotoneaster, Crataegus, Dryas, Exochorda, Filipendula,
Geum, Gillenia, Holodiscus, Ivesia, Malus, Petrophytum, Photinia, Potentilla,
Prinsepia, Prunus, Pyracantha, Pyrus, Rhodotypos, Rosa, Rubus, Sanguisorba,
Sibbaldia, Sorbus, Spiraea, Stephanandra, Waldsteinia: Low temperature and
multicycle germinations were found.
\item[Rubiaceae] Asperula, Cephalanthus, Cruckshanksia, Hedyotis, Houstonia,
Mitchella, Mussaenda: An absolute photorequirement is found in several genera.
\item[Rutaceae] Citrus, Crowea, Dictamnus, Evodia, Phellodendron, Poncirus,
Ptelea, Skimmia: Germinations are sometimes extended in multicycle patterns.
\item[Salicaceae] Populus, Salix
\item[Sapindaceae] Cardiospermum, Koelreuteria
\item[Sarraceniaceae] Sarracenia
\item[Saxifragaceae] Astelboides, Astilbe, Bensoniella, Bergenia, Deutzia,
Elmera, Heuchera, Hydrangea, Itea, Kirengeshoma, Leptarrhena, Mitella, Parnassia,
Peltoboykinia, Philadelphus, Ribes, Rodgersia, Saxifraga, Telesonix, Tellima, Tiarella:
Photorequirements were frequent.
\item[Schisandraceae] Schisandra
\item[Scrophulariaceae] Agalinis, Alonsoa, Antirrhinum, Asarina,
Aureolaria, Calceolaria, Castilleja, Chaenorrhinum, Chelone, Digitalis, Diplacus,
Erinus, Gerardia, Hebe, Lagotis, Libanotis, Linaria, Melampyrum, Mimulus, Nemesia,
Ourisia, Paulownia, Pedicularis, Penstemon, Striga, Synthyris, Verbascum, Veronica,
Wulfenia, Zaluzianskya: Photorequirements were frequent and sometimes complex.
\item[Simaroubaceae] Ailanthus
\item[Solanaceae] Atropa, Capsicum, Datura, Lycopersicon, Petunia, Physalis,
Salpiglosis, Schizanthus, Solanum: Photorequirements and stimulation by GA-3
were frequent.
\item[Styracaceae] Halesia, Styrax: These are truly recalcitrant germinators.
Although all species gave some germination, the critical factors are not yet
understood.
\item[Svmplocaceae] Symplocos
\item[Taxaceae] Taxus: Germinations were extended.
\item[Tecophiliaceae] Conanthera
\item[Theaceae] Franklinia, Stewartia
\item[Thymelaeaceae] Daphne, Pimelea, Stellaria: Germinations were
extended.
\item[Tiliaceae] Tilia: Germinations were extended.
\item[Tropaeolaceae] Tropaeolum
\item[Typhaceae] Typha
\item[Ulmaceae] Aphananthe, Celtis, Ulmus
\item[Valerianaceae] Centranthus
\item[Verbenaceae] Callicarpa, Caryopteris, Pleuroginella, Verbena:
Photorequirements were found and they were sometimes complex.
\item[Violaceae] Viola: Photorequirements and stimulation by GA-3 were
frequent.
\item[Vitaceae] Parthenocissus, Vitis
\item[Zygophyllaceae] Larrea
% \item[Name in doubt: Grumolo]
\end{description}

